\chapter[From Bloch Waves to Effective-Mass Theory]{From Bloch Waves to
Effective-Mass Theory: A Controlled Atomistic-to-Continuum Framework for
Semiconductor Physics}
\label{ch:operator-reduction-motivation}

\textbf{Status: conceptual and illustrative framework. The central hypothesis,
continuum criteria, and intellectual contributions are proposed work; the
particle-in-a-box example is not research evidence or a validated reduction of
silicon.}

\section{Project motivation and significance}

Semiconductor physics is commonly taught through two quantum-mechanical
descriptions that are physically related but mathematically disconnected.
Students first solve a microscopic periodic Hamiltonian,
\begin{equation}
    \hat H_{\mathrm{micro}}
    =
    -\frac{\hbar^2}{2m_0}\nabla^2
    +
    V_{\mathrm{per}}(\mathbf r),
\end{equation}
and obtain Bloch states and energy bands. They are later introduced to
effective-mass theory through a continuum envelope Hamiltonian. Written
relative to the selected band-edge energy, it is
\begin{equation}
    \hat H_{\mathrm{EMT}}
    =
    -\frac{\hbar^2}{2}
    \nabla\mathbin{\cdot}\mathbf m^{*-1}\nabla
    +
    V_{\mathrm{ext}}(\mathbf r).
\end{equation}

For a smooth, nondegenerate band near an extremum $\mathbf k_0$, the usual
connection expands the dispersion as
\begin{equation}
    E_n(\mathbf k_0+\mathbf q)
    =
    E_n(\mathbf k_0)
    +
    \frac{\hbar^2}{2}
    \mathbf q^{\mathsf T}\mathbf m_n^{*-1}\mathbf q
    +
    O\!\left(\lVert\mathbf q\rVert^3\right).
\end{equation}
After subtracting the constant $E_n(\mathbf k_0)\hat I$ and replacing
$\mathbf q$ by $-i\nabla$, this argument explains the form of the effective
kinetic operator. It does not construct the envelope-function state space,
define the map from microscopic Bloch states to continuum states, or quantify
the approximations introduced by band truncation, gauge selection, valley
decomposition, and spatial coarse-graining.

Effective-mass theory is therefore often presented as a second quantum
mechanics rather than as a controlled reduction of the first. The
effective-mass tensor is obtained from the curvature of a Bloch band, but the
corresponding state-space reduction is usually left implicit.

This project began from a practical attempt to close that pedagogical gap. The
initial objective was to construct an explicit periodic potential that students
could discretize and diagonalize, yet which remained sufficiently expressive
to produce an indirect band gap. Such a model would provide an executable
bridge between elementary periodic quantum mechanics and the band structure of
an indirect-gap semiconductor.

The most direct strategy appears elementary: calculate Kohn--Sham
density-functional-theory Bloch eigenpairs, reconstruct a retained Hamiltonian,
subtract the known kinetic operator, and interpret the remainder as an
effective periodic potential.

Let $\hat H_{\mathrm{KS}}(\mathbf k)$ be a Kohn--Sham Bloch-fiber operator
with cell-periodic eigenstates $\lvert u_{n\mathbf k}\rangle$. For a retained
set of bands $\mathcal B$, define
\begin{equation}
    \hat P(\mathbf k)
    =
    \sum_{n\in\mathcal B}
    \lvert u_{n\mathbf k}\rangle
    \langle u_{n\mathbf k}\rvert
\end{equation}
and
\begin{equation}
    \hat H_{\mathrm{KS}}^{(P)}(\mathbf k)
    =
    \left.
        \hat P(\mathbf k)
        \hat H_{\mathrm{KS}}(\mathbf k)
        \hat P(\mathbf k)
    \right|_{\mathcal H^{(P)}(\mathbf k)}
    =
    \sum_{n\in\mathcal B}
    \varepsilon_n(\mathbf k)
    \lvert u_{n\mathbf k}\rangle
    \langle u_{n\mathbf k}\rvert,
\end{equation}
where
\begin{equation}
    \mathcal H^{(P)}(\mathbf k)
    =
    \operatorname{Ran}\hat P(\mathbf k).
\end{equation}
For entangled bands, constructing a smooth retained family may require
projection or disentanglement rather than selection by a single global energy
cutoff.

The apparently obvious construction is
\begin{equation}
    \hat V_{\mathrm{eff}}^{(P)}(\mathbf k)
    \stackrel{?}{=}
    \hat H_{\mathrm{KS}}^{(P)}(\mathbf k)
    -
    \hat T^{(P)}(\mathbf k),
\end{equation}
where
\begin{equation}
    \hat T^{(P)}(\mathbf k)
    =
    \left.
        \hat P(\mathbf k)
        \hat T(\mathbf k)
        \hat P(\mathbf k)
    \right|_{\mathcal H^{(P)}(\mathbf k)}.
\end{equation}
This subtraction is mathematically valid on the retained fiber, but its result
is a projected operator, not necessarily a unique local scalar potential.
Projection may induce nonlocality, while nonlocal pseudopotentials, orbital
structure, spin--orbit interactions, and band truncation may introduce
components that cannot be represented by a simple function $V(\mathbf r)$.
Bloch gauge freedom further changes the matrix coordinates of the retained
operator without changing its abstract action or spectrum.

The desired teaching potential must therefore be obtained by reducing the
operator residual into an explicitly chosen model class:
\begin{equation}
    \hat V_{\mathrm{teach}}
    =
    \operatorname*{arg\,min}_{\hat V\in\mathcal V_{\mathrm{adm}}}
    \sum_{\mathbf k\in\mathcal K}
    \left\|
        \hat H_{\mathrm{KS}}^{(P)}(\mathbf k)
        -
        \hat T^{(P)}(\mathbf k)
        -
        \left.
            \hat P(\mathbf k)
            \hat V(\mathbf k)
            \hat P(\mathbf k)
        \right|_{\mathcal H^{(P)}(\mathbf k)}
    \right\|_*^2.
\end{equation}
Here, $\mathcal K$ is the sampled reciprocal-space region and
$\mathcal V_{\mathrm{adm}}$ may represent local periodic potentials, truncated
Fourier expansions, finite real-space stencils, or constrained tight-binding
operator classes. The norm $\lVert\cdot\rVert_*$ and any reciprocal-space
weights must be chosen to measure operator error consistently with the physical
observables and continuum limit of interest. This is a proposed model-class
problem, not a presently accepted fitting criterion.

The finite example developed next separates exact operator retention from
model-class approximation and shows why all components of an additive
Hamiltonian must be reduced through the same map. The later sections then
return to pristine--doped comparison, continuum admissibility, and the broader
research program.

\section{From spectral data to an effective Hamiltonian}

The construction in this section is a pedagogical discrete-spectrum case. It
is not the literal subspace construction used for a periodic crystal, where
Bloch-fiber representations and possibly projection or disentanglement enter.

The genesis of this problem begins with the spectral problem for a quantum
Hamiltonian $\hat H$,
\begin{equation}
    \hat H\lvert\psi_i\rangle
    =
    E_i\lvert\psi_i\rangle,
\end{equation}
where $(E_i,\lvert\psi_i\rangle)$ is the $i$th eigenpair. Assuming that the
relevant spectrum is discrete and that the eigenstates form a complete
orthonormal basis of the Hilbert space $\mathcal H$, the Hamiltonian admits the
spectral decomposition
\begin{equation}
    \hat H
    =
    \sum_i
    E_i\lvert\psi_i\rangle\langle\psi_i\rvert.
\end{equation}

Suppose that the physical questions of interest involve only states below an
energy cutoff $E_{\mathrm{cut}}$. We define the retained low-energy subspace
\begin{equation}
    \mathcal H^{(P)}
    =
    \operatorname{span}
    \left\{
        \lvert\psi_i\rangle
        \;\middle|\;
        E_i\leq E_{\mathrm{cut}}
    \right\},
\end{equation}
with orthogonal projector
\begin{equation}
    \hat P
    =
    \sum_{E_i\leq E_{\mathrm{cut}}}
    \lvert\psi_i\rangle\langle\psi_i\rvert.
\end{equation}

The corresponding retained Hamiltonian is
\begin{equation}
    \hat H^{(P)}
    =
    \left.
        \hat P\hat H\hat P
    \right|_{\mathcal H^{(P)}}
    =
    \sum_{E_i\leq E_{\mathrm{cut}}}
    E_i\lvert\psi_i\rangle\langle\psi_i\rvert.
\end{equation}
It exactly reproduces every retained eigenpair,
\begin{equation}
    \hat H^{(P)}\lvert\psi_i\rangle
    =
    E_i\lvert\psi_i\rangle,
    \qquad
    E_i\leq E_{\mathrm{cut}},
\end{equation}
but it does not specify the action of $\hat H$ outside
$\mathcal H^{(P)}$. It is therefore an exact restriction of the parent
Hamiltonian, not yet an approximate effective model.

We reserve the notation $\hat H_{\mathrm{eff}}(\boldsymbol\theta)$ for an
operator belonging to a prescribed effective-model class with parameters
$\boldsymbol\theta$. The distinction is important:
\begin{equation}
    \underbrace{\hat H^{(P)}}_{\text{exact retained operator}}
    \qquad\text{and}\qquad
    \underbrace{\hat H_{\mathrm{eff}}(\boldsymbol\theta)}_{
        \text{approximate model on an identified retained space}}
\end{equation}
are not the same construction.

\section{Computational example: a particle in a one-dimensional box}

The particle in a one-dimensional box provides the simplest computational
setting in which the continuous operator, its discrete matrix representation,
its retained subspace, and an effective model can be separated explicitly.

\subsection{Continuous problem}

Consider a particle of mass $m$ confined to the interval
\begin{equation}
    \Omega=(0,L).
\end{equation}
Inside the box, the Hamiltonian is
\begin{equation}
    \hat H
    =
    -\frac{\hbar^2}{2m}\frac{d^2}{dx^2},
\end{equation}
with operator domain
\begin{equation}
    \mathcal D(\hat H)
    =
    H^2(0,L)\cap H_0^1(0,L).
\end{equation}
Equivalently, the wavefunction satisfies the Dirichlet boundary conditions
\begin{equation}
    \psi(0)=\psi(L)=0.
\end{equation}

The exact eigenfunctions and eigenvalues are
\begin{align}
    \psi_n(x)
    &=
    \sqrt{\frac{2}{L}}
    \sin\left(\frac{n\pi x}{L}\right),
    \\
    E_n
    &=
    \frac{\hbar^2\pi^2n^2}{2mL^2},
    \qquad n=1,2,\ldots.
\end{align}

\subsection{Finite-difference representation}

Introduce $N$ interior grid points
\begin{equation}
    x_j=j\Delta x,
    \qquad
    j=1,\ldots,N,
\end{equation}
where
\begin{equation}
    \Delta x=\frac{L}{N+1}.
\end{equation}
The boundary points are excluded from the numerical state vector because the
Dirichlet conditions fix their values to zero. A sampled wavefunction is
represented by
\begin{equation}
    \boldsymbol\psi
    =
    \begin{bmatrix}
        \psi(x_1) & \psi(x_2) & \cdots & \psi(x_N)
    \end{bmatrix}^{\mathsf T}
    \in\mathbb C^N.
\end{equation}

Using the centered finite-difference approximation
\begin{equation}
    \left.
        \frac{d^2\psi}{dx^2}
    \right|_{x=x_j}
    \approx
    \frac{
        \psi(x_{j+1})-2\psi(x_j)+\psi(x_{j-1})
    }{\Delta x^2},
\end{equation}
the Hamiltonian is represented by the Hermitian matrix
\begin{equation}
    \mathbf H
    =
    \frac{\hbar^2}{2m\Delta x^2}
    \begin{bmatrix}
        2  & -1 & 0  & \cdots & 0 \\
        -1 & 2  & -1 & \ddots & \vdots \\
        0  & -1 & 2  & \ddots & 0 \\
        \vdots & \ddots & \ddots & \ddots & -1 \\
        0 & \cdots & 0 & -1 & 2
    \end{bmatrix}.
\end{equation}

The discrete eigenvalue problem is
\begin{equation}
    \mathbf H\mathbf v_i
    =
    \varepsilon_i\mathbf v_i,
\end{equation}
where $\varepsilon_i$ approximates $E_i$. With Euclidean normalization,
$\mathbf v_i$ approximates the quadrature-scaled samples
$\sqrt{\Delta x}\,\psi_i(x_j)$ rather than the unscaled samples themselves.
The finite-difference eigenvalues are
\begin{equation}
    \varepsilon_i
    =
    \frac{2\hbar^2}{m\Delta x^2}
    \sin^2\left(
        \frac{i\pi}{2(N+1)}
    \right),
    \qquad i=1,\ldots,N.
\end{equation}
For fixed $i$,
\begin{equation}
    \lim_{\Delta x\rightarrow 0}\varepsilon_i
    =
    E_i.
\end{equation}
The low-energy eigenpairs converge first because their wavelengths remain
large compared with the grid spacing. States near the grid-resolution limit
are strongly affected by the finite-difference dispersion relation.

\subsection{Spectral reconstruction in the discrete space}

Let
\begin{equation}
    \boldsymbol\Phi
    =
    \begin{bmatrix}
        \mathbf v_1 & \mathbf v_2 & \cdots & \mathbf v_N
    \end{bmatrix}
\end{equation}
be the unitary eigenvector matrix and let
\begin{equation}
    \boldsymbol\Lambda
    =
    \operatorname{diag}
    \left(
        \varepsilon_1,\ldots,\varepsilon_N
    \right).
\end{equation}
Then
\begin{equation}
    \boldsymbol\Phi^\dagger\boldsymbol\Phi
    =
    \mathbf I,
\end{equation}
and the discrete Hamiltonian can be reconstructed exactly as
\begin{equation}
    \mathbf H
    =
    \boldsymbol\Phi
    \boldsymbol\Lambda
    \boldsymbol\Phi^\dagger
    =
    \sum_{i=1}^{N}
    \varepsilon_i\mathbf v_i\mathbf v_i^\dagger.
\end{equation}

The relation between the continuous and discrete objects should therefore be
written as a representation map,
\begin{equation}
    \hat H
    \longrightarrow
    \mathbf H,
\end{equation}
rather than as a literal equality. The matrix $\mathbf H$ depends on the
chosen grid, finite-difference stencil, boundary conditions, and units.

\subsection{Low-energy retention}

Assume $1\leq M<N$. Let
\begin{equation}
    \boldsymbol\Phi_M
    =
    \begin{bmatrix}
        \mathbf v_1 & \mathbf v_2 & \cdots & \mathbf v_M
    \end{bmatrix}
    \in\mathbb C^{N\times M}
\end{equation}
contain the lowest $M$ eigenvectors, and define
\begin{equation}
    \boldsymbol\Lambda_M
    =
    \operatorname{diag}
    \left(
        \varepsilon_1,\ldots,\varepsilon_M
    \right).
\end{equation}
The orthogonal projector onto the retained discrete subspace is
\begin{equation}
    \mathbf P_M
    =
    \boldsymbol\Phi_M\boldsymbol\Phi_M^\dagger.
\end{equation}

The retained Hamiltonian embedded in the original grid space is
\begin{align}
    \mathbf H^{(P_M)}
    &=
    \mathbf P_M\mathbf H\mathbf P_M
    \\
    &=
    \boldsymbol\Phi_M
    \boldsymbol\Lambda_M
    \boldsymbol\Phi_M^\dagger.
\end{align}
This is an $N\times N$ matrix of rank $M$. It represents the exact action of
$\mathbf H$ on $\operatorname{Ran}(\mathbf P_M)$ and vanishes on the
orthogonal complement.

The same retained operator has the $M\times M$ matrix representation
\begin{align}
    \mathbf H_{\mathrm{red}}
    &=
    \boldsymbol\Phi_M^\dagger
    \mathbf H
    \boldsymbol\Phi_M
    \\
    &=
    \boldsymbol\Lambda_M
\end{align}
in the retained eigenbasis. Thus,
\begin{equation}
    \mathbf H^{(P_M)}\in\mathbb C^{N\times N},
    \qquad
    \mathbf H_{\mathrm{red}}\in\mathbb C^{M\times M},
\end{equation}
The two matrices represent the same retained action only after restricting
$\mathbf H^{(P_M)}$ to $\operatorname{Ran}(\mathbf P_M)$ and identifying
$\mathbb C^M$ with that range through $\boldsymbol\Phi_M$.

Consequently, the expression
\begin{equation}
    \mathbf H-\mathbf H_{\mathrm{red}}
\end{equation}
is undefined when $M\neq N$. The reduced matrix must first be embedded into
the grid space,
\begin{equation}
    \mathbf H^{(P_M)}
    =
    \boldsymbol\Phi_M
    \mathbf H_{\mathrm{red}}
    \boldsymbol\Phi_M^\dagger,
\end{equation}
or the full Hamiltonian must be restricted to the retained coordinates.

\subsection{Operator residuals}

The full-grid residual associated with spectral retention is
\begin{align}
    \mathbf R_{\mathrm{full}}
    &=
    \mathbf H-\mathbf H^{(P_M)}
    \\
    &=
    \sum_{i=M+1}^{N}
    \varepsilon_i\mathbf v_i\mathbf v_i^\dagger.
\end{align}
For the positive, increasingly ordered particle-in-a-box spectrum,
\begin{align}
    \left\|\mathbf R_{\mathrm{full}}\right\|_2
    &=
    \varepsilon_N,
    \\
    \left\|\mathbf R_{\mathrm{full}}\right\|_{\mathrm F}
    &=
    \left(
        \sum_{i=M+1}^{N}\varepsilon_i^2
    \right)^{1/2}.
\end{align}
These norms are dominated by the discarded high-energy sector. They do not
measure the accuracy with which the retained low-energy action is reproduced.
Indeed,
\begin{equation}
    \mathbf P_M
    \left(
        \mathbf H-\mathbf H^{(P_M)}
    \right)
    \mathbf P_M
    =
    \mathbf 0.
\end{equation}

Now suppose that
$\mathbf H_{\mathrm{eff}}(\boldsymbol\theta)\in\mathbb C^{M\times M}$ is a
parameterized effective model expressed in the retained basis. Its relevant
operator residual is
\begin{equation}
    \mathbf R_H^{(P_M)}(\boldsymbol\theta)
    =
    \boldsymbol\Phi_M^\dagger
    \mathbf H
    \boldsymbol\Phi_M
    -
    \mathbf H_{\mathrm{eff}}(\boldsymbol\theta).
\end{equation}
A Frobenius-norm fitting problem can then be posed as
\begin{equation}
    \boldsymbol\theta^\star
    =
    \operatorname*{arg\,min}_{\boldsymbol\theta}
    \left\|
        \mathbf R_H^{(P_M)}(\boldsymbol\theta)
    \right\|_{\mathrm F}^2.
\end{equation}
This comparison is well defined because both matrices act on the same
$M$-dimensional state space and use the same retained basis.

\section{Conditions for identifying a reduced potential operator}

A reduced Hamiltonian does not by itself determine a unique kinetic--potential
decomposition. Such an identification requires a prescribed model class and a
specified kinetic operator; projected components may also be nonlocal even
when their parent-space counterparts are local.

Suppose that the parent Hamiltonian admits the additive decomposition
\begin{equation}
    \hat H=\hat T+\hat V,
\end{equation}
where $\hat T$ is the kinetic-energy operator and $\hat V$ is the
potential-energy operator. In the parent state space,
\begin{equation}
    \hat V=\hat H-\hat T.
\end{equation}
After discretization on a common numerical state space, this becomes
\begin{align}
    \mathbf H
    &=
    \mathbf T+\mathbf V,
    \\
    \mathbf V
    &=
    \mathbf H-\mathbf T.
\end{align}

It is tempting to apply the same construction to a reduced Hamiltonian and
write
\begin{equation}
    \mathbf V_{\mathrm{eff}}
    \stackrel{?}{=}
    \mathbf H_{\mathrm{eff}}-\mathbf T.
\end{equation}
This expression is meaningful only if $\mathbf H_{\mathrm{eff}}$ and
$\mathbf T$ act on the same state space and are represented in the same
basis. If the Hamiltonian has been reduced, then the kinetic operator must be
reduced through the same map.

\subsection{Consistent reduction of an additive decomposition}

Define the projection map
\begin{equation}
    \mathcal R_{P_M}(\mathbf A)
    =
    \mathbf P_M\mathbf A\mathbf P_M.
\end{equation}
Because $\mathcal R_{P_M}$ is linear,
\begin{align}
    \mathcal R_{P_M}(\mathbf H)
    &=
    \mathcal R_{P_M}(\mathbf T)
    +
    \mathcal R_{P_M}(\mathbf V),
    \\
    \mathbf H^{(P_M)}
    &=
    \mathbf T^{(P_M)}
    +
    \mathbf V^{(P_M)},
\end{align}
where
\begin{align}
    \mathbf T^{(P_M)}
    &=
    \mathbf P_M\mathbf T\mathbf P_M,
    \\
    \mathbf V^{(P_M)}
    &=
    \mathbf P_M\mathbf V\mathbf P_M.
\end{align}
The retained potential is therefore
\begin{equation}
    \mathbf V^{(P_M)}
    =
    \mathbf H^{(P_M)}
    -
    \mathbf T^{(P_M)},
\end{equation}
not
\begin{equation}
    \mathbf V^{(P_M)}
    \stackrel{?}{=}
    \mathbf H^{(P_M)}-\mathbf T.
\end{equation}

In the retained basis, the same relation is
\begin{align}
    \mathbf H_{\mathrm{red}}
    &=
    \boldsymbol\Phi_M^\dagger
    \mathbf H
    \boldsymbol\Phi_M,
    \\
    \mathbf T_{\mathrm{red}}
    &=
    \boldsymbol\Phi_M^\dagger
    \mathbf T
    \boldsymbol\Phi_M,
    \\
    \mathbf V_{\mathrm{red}}
    &=
    \boldsymbol\Phi_M^\dagger
    \mathbf V
    \boldsymbol\Phi_M,
\end{align}
with
\begin{equation}
    \mathbf V_{\mathrm{red}}
    =
    \mathbf H_{\mathrm{red}}-\mathbf T_{\mathrm{red}}.
\end{equation}

\subsection{The particle-in-a-box counterexample}

For the infinite particle in a box, the interior potential vanishes:
\begin{equation}
    V(x)=0.
\end{equation}
The infinite walls are not represented by finite values of $V(x)$; they are
encoded in the domain of $\hat H$ and its Dirichlet boundary conditions. On
the interior finite-difference grid,
\begin{equation}
    \mathbf H=\mathbf T,
    \qquad
    \mathbf V=\mathbf 0.
\end{equation}
Consistent projection gives
\begin{equation}
    \mathbf H^{(P_M)}
    =
    \mathbf T^{(P_M)},
\end{equation}
and hence
\begin{equation}
    \mathbf V^{(P_M)}
    =
    \mathbf H^{(P_M)}-\mathbf T^{(P_M)}
    =
    \mathbf 0.
\end{equation}

Now subtract the unprojected kinetic matrix instead:
\begin{equation}
    \widetilde{\mathbf V}_{\mathrm{eff}}
    =
    \mathbf H^{(P_M)}-\mathbf T.
\end{equation}
For the particle in the box,
\begin{equation}
    \widetilde{\mathbf V}_{\mathrm{eff}}
    =
    \mathbf P_M\mathbf T\mathbf P_M-\mathbf T.
\end{equation}
Because $\mathbf P_M$ is a spectral projector of $\mathbf T$, it commutes
with $\mathbf T$. Let
\begin{equation}
    \mathbf Q_M=\mathbf I-\mathbf P_M
\end{equation}
project onto the discarded subspace. Then
\begin{align}
    \widetilde{\mathbf V}_{\mathrm{eff}}
    &=
    -\mathbf Q_M\mathbf T\mathbf Q_M
    \\
    &=
    -\sum_{i=M+1}^{N}
    \varepsilon_i\mathbf v_i\mathbf v_i^\dagger.
\end{align}
The reconstructed matrix is nonzero even though the physical potential inside
the box is zero. It is not an effective potential. It is the negative of the
discarded high-energy kinetic operator produced by combining a projected
Hamiltonian with an unprojected kinetic operator.

This counterexample separates two obstructions. Subtracting
$\mathbf H_{\mathrm{red}}$ from the full-grid $\mathbf T$ is undefined because
the matrices act on different spaces. Subtracting $\mathbf H^{(P_M)}$ from
$\mathbf T$ is algebraically defined in the common grid representation, but
the result is not the consistently reduced potential. Physical interpretation
therefore requires both a common representation and a common reduction map.

\subsection{Potential residual}

Let $\mathbf V_{\mathrm{eff}}(\boldsymbol\theta)$ be a proposed effective
potential acting on the retained space. The potential residual is
\begin{equation}
    \mathbf R_V^{(P_M)}(\boldsymbol\theta)
    =
    \mathbf V_{\mathrm{red}}
    -
    \mathbf V_{\mathrm{eff}}(\boldsymbol\theta).
\end{equation}
If
\begin{align}
    \mathbf V_{\mathrm{red}}
    &=
    \mathbf H_{\mathrm{red}}-\mathbf T_{\mathrm{red}},
    \\
    \mathbf V_{\mathrm{eff}}(\boldsymbol\theta)
    &=
    \mathbf H_{\mathrm{eff}}(\boldsymbol\theta)
    -
    \mathbf T_{\mathrm{eff}}(\boldsymbol\theta),
\end{align}
then
\begin{equation}
    \mathbf R_V^{(P_M)}(\boldsymbol\theta)
    =
    \left(
        \mathbf H_{\mathrm{red}}
        -
        \mathbf H_{\mathrm{eff}}(\boldsymbol\theta)
    \right)
    -
    \left(
        \mathbf T_{\mathrm{red}}
        -
        \mathbf T_{\mathrm{eff}}(\boldsymbol\theta)
    \right).
\end{equation}
Only when the kinetic operators have been identified,
\begin{equation}
    \mathbf T_{\mathrm{eff}}(\boldsymbol\theta)
    =
    \mathbf T_{\mathrm{red}},
\end{equation}
does the potential residual reduce to the Hamiltonian residual:
\begin{equation}
    \mathbf R_V^{(P_M)}(\boldsymbol\theta)
    =
    \mathbf H_{\mathrm{red}}
    -
    \mathbf H_{\mathrm{eff}}(\boldsymbol\theta).
\end{equation}

\section{From the particle in a box to electronic-structure reduction}

The particle-in-a-box example is elementary, but the obstruction it reveals is
general. In a common parent state space, the decomposition
\begin{equation}
    \mathbf H=\mathbf T+\mathbf V
\end{equation}
is meaningful because all three matrices act on the same discrete state space
and are represented in the same basis. Reduction preserves this additive
decomposition only when every term is reduced through the same map:
\begin{equation}
    \mathcal R_{P_M}(\mathbf H-\mathbf T)
    =
    \mathcal R_{P_M}(\mathbf H)
    -
    \mathcal R_{P_M}(\mathbf T).
\end{equation}
In general,
\begin{equation}
    \mathcal R_{P_M}(\mathbf H-\mathbf T)
    \neq
    \mathcal R_{P_M}(\mathbf H)-\mathbf T.
\end{equation}

The distinction becomes more consequential in electronic-structure
calculations. A Kohn--Sham Hamiltonian, a Wannier Hamiltonian, a tight-binding
Hamiltonian, and an effective-mass Hamiltonian generally do not begin as
matrices acting on an identical numerical state space. They may differ in
dimension, basis, gauge, spatial resolution, boundary conditions, geometry,
units, and energy reference. Their matrices therefore cannot be physically
compared merely because they are all called Hamiltonians.

The central problem of this work is consequently not the mechanical
subtraction of two matrices. It is the construction and validation of the
state-space identifications that make such a subtraction mathematically
defined and physically interpretable.

\section{From periodic potentials to impurity operators}

The same structural problem becomes more consequential when the model is
extended from bulk bands to semiconductor impurities. Given pristine and doped
Hamiltonians, the natural proposal is
\begin{equation}
    \hat V_{\mathrm{imp}}
    \stackrel{?}{=}
    \hat H_d-\hat H_b.
\end{equation}
This expression is not yet an operator identity because independently
constructed pristine and doped calculations may use different retained
subspaces, Bloch or Wannier gauges, basis orderings, geometries, spin
structures, energy references, and numerical dimensions. Their matrices cannot
be interpreted as representations in common physical coordinates without an
explicit identification.

Let $s\in\{b,d\}$ denote matched bulk and doped branches, with retained
operators
\begin{equation}
    \hat H_s^{(P)}
    =
    \left.
        \hat P_s\hat H_s\hat P_s
    \right|_{\mathcal H_s^{(P)}},
    \qquad
    \mathcal H_s^{(P)}
    =
    \operatorname{Ran}\hat P_s.
\end{equation}
If the retained ranks agree, a meaningful comparison requires an
identification
\begin{equation}
    \hat U_d:
    \mathcal H_b^{(P)}
    \longrightarrow
    \mathcal H_d^{(P)}.
\end{equation}
Let $c_b$ and $c_d$ be the scalar reference energies selected by the declared
alignment procedure, and define
\begin{equation}
    \overline{\hat H}_s^{(P)}
    =
    \hat H_s^{(P)}-c_s\hat I_s.
\end{equation}
For a unitary identification $\hat U_d$, the aligned impurity operator on the
bulk retained space is
\begin{equation}
    \Delta\hat H_d^{(P)}
    =
    \hat U_d^\dagger
    \overline{\hat H}_d^{(P)}
    \hat U_d
    -
    \overline{\hat H}_b^{(P)}.
\end{equation}
Unequal retained ranks require a separately defined partial identification or
further model reduction; they do not admit this unitary formula unchanged.

Equal-dimensional spaces admit many abstract unitary maps. Most are not
physically admissible: an arbitrary unitary can scramble localization, orbital
character, crystal momentum, valley content, symmetry, and spin while leaving
the spectrum unchanged. Gauge alignment is therefore necessary to define the
subtraction, but it is not sufficient to select a physical impurity operator.
The identification must also respect the structures needed by the intended
reduced theory.

\section{Central hypothesis}

The central hypothesis proposed by this project is:
\begin{quote}
\itshape
The intended continuum limit is not merely a downstream approximation to an
already constructed atomistic model. It constrains, from the beginning, the
admissible retained spaces, gauges, identification maps, and reduced operator
classes.
\end{quote}

For a semiconductor representation to support an effective-mass
interpretation, its atomistic-to-continuum identification must preserve or
control the structures used by the continuum theory. These include
\begin{itemize}
    \item real-space localization and lattice translation;
    \item Bloch wave vector and valley correspondence;
    \item orbital, symmetry, and spin content;
    \item smoothness across the relevant reciprocal-space region;
    \item common bulk behavior away from an impurity;
    \item boundary conditions and normalization; and
    \item separation of slowly varying envelopes from lattice-periodic
          components.
\end{itemize}

Let
\begin{equation}
    \epsilon=\frac{a}{L}
\end{equation}
denote the ratio of lattice scale $a$ to an envelope or impurity scale $L$.
For a specified family of retained atomistic operators
$\hat H_\epsilon^{(P)}$ and appropriately normalized identification maps
\begin{equation}
    \hat J_\epsilon:
    \mathcal H_\epsilon^{(P)}
    \longrightarrow
    \mathcal H_{\mathrm{EMT}},
\end{equation}
one possible continuum criterion is norm-resolvent convergence,
\begin{equation}
    \left\|
        \hat J_\epsilon
        \left(\hat H_\epsilon^{(P)}-z\right)^{-1}
        \hat J_\epsilon^\dagger
        -
        \left(\hat H_{\mathrm{EMT}}-z\right)^{-1}
    \right\|
    \longrightarrow 0,
    \qquad
    \epsilon\longrightarrow 0,
\end{equation}
for $z$ in an identified common resolvent set. This is a candidate analytical
form, not an established result of the project. A later specification may
instead require a weaker, explicitly defined combination of spectral,
subspace, operator, and observable convergence criteria.

The continuum limit thereby becomes an upstream admissibility condition. A
gauge, basis, identification, or reduced operator class is acceptable only if
it preserves, or approximates with controlled error, the microscopic structures
needed to recover the declared envelope theory.

\section{Relation to Hilbert's sixth problem}

This research is situated within the broad mathematical tradition associated
with Hilbert's sixth problem: the rigorous formulation of physical theories and
the justification of limiting relations between microscopic and continuum
descriptions. It does not claim to solve Hilbert's sixth problem in its general
form. Instead, it addresses a bounded and experimentally consequential
instance: the reduction of microscopic electronic-structure operators to
continuum envelope models for indirect-gap semiconductors and their impurities.

The relevant limiting process is not only a passage between different length
scales. It is a passage between different state spaces, operator algebras,
coordinate systems, and model classes:
\begin{equation}
    \text{KS--DFT}
    \longrightarrow
    \text{retained Bloch operator}
    \longrightarrow
    \text{localized atomistic operator}
    \longrightarrow
    \text{effective-mass continuum operator}.
\end{equation}
Each arrow requires an explicit map, a statement of preserved structure, and an
error criterion. Without these elements, agreement between selected band
energies does not establish that two models represent the same reduced physics.

The effective-mass limit is therefore both an endpoint and a constraint. The
program asks not only whether an atomistic calculation can reproduce a
continuum observable, but which atomistic representations and reductions are
capable of supporting the continuum theory at all.

\section{Intellectual contribution}

The proposed program aims to develop a gauge-equivariant,
continuum-constrained framework for semiconductor operator reduction. Its
intended contributions are
\begin{enumerate}
    \item reconstruction of retained operators from Kohn--Sham Bloch
          eigenpairs;
    \item separation of operator reconstruction from model-class
          approximation;
    \item quantitative projection onto interpretable periodic-potential,
          stencil, and tight-binding classes;
    \item explicit alignment of pristine and doped retained state spaces;
    \item construction of impurity operators whose abstract action is invariant
          under compatible coordinate changes after alignment;
    \item formulation of continuum-compatibility criteria using operator,
          subspace, spectral, and observable errors; and
    \item an executable pedagogical chain connecting periodic potentials,
          indirect band structure, impurity operators, and effective-mass
          theory.
\end{enumerate}

The conceptual distinction is concise:
\begin{quote}
\itshape
Gauge alignment makes an operator comparison mathematically defined;
continuum compatibility makes the resulting reduced model physically
admissible.
\end{quote}

What began as an attempt to construct a teachable periodic potential therefore
exposes a broader mathematical problem. Obtaining the potential requires
reconstructing an operator rather than merely fitting a spectrum. Introducing
an impurity requires comparing independently constructed operators. Comparing
those operators requires identifying their state spaces and gauges. Selecting
physically admissible identifications and reduced models requires a declared
continuum limit.

This chain provides both the scientific rationale and the educational
significance of the proposed research. It converts the normally implicit
transition from Bloch theory to effective-mass theory into an explicit,
testable, and executable sequence of operator reductions while keeping every
approximation and evidence claim explicit. Chapter~\ref{ch:physical-problem}
next specializes this framework to the silicon host and its substitutional
phosphorus and boron branches.
